\documentclass[11pt]{article}
\usepackage[margin=2cm]{geometry}
\usepackage{booktabs}
\usepackage{multirow}
\usepackage{array}
\usepackage[table]{xcolor}
\usepackage{amsmath}

\definecolor{best}{RGB}{0,90,200}
\newcommand{\bb}[1]{\textcolor{best}{\textbf{#1}}}

\begin{document}

\begin{center}
{\Large\textbf{Bitcoin Forecasting --- Experiment Results}}
\end{center}

\medskip
All evaluations use walk-forward folds, $k{=}6$ horizon, and 0.1\% round-trip
slippage. Our models use a frozen Qwen3-VL-4B backbone with a DiT flow-matching
head; \textbf{v15} adds per-asset target normalisation. \textbf{Kronos}
(NeoQuasar/Kronos-base) is a pretrained K-line foundation model, reported only on
fold6 (2025), the single window after its pretraining cutoff (June 2024).

%==========================================================================
\section*{Part A --- Crypto experiments (24/7, 6 folds 2020--2025)}
%==========================================================================

% ----- Table A1 -----
\begin{table}[h]\centering
\caption{Directional accuracy and regime bias. \textbf{Long\_ns std} (across the
6 folds) measures regime memorisation --- lower means less polarised toward one
direction per year. Normalisation (v15) roughly halves it.}
\label{tab:dir}
\begin{tabular}{llccc}
\toprule
Model & Period & DA & Long\_ns & Long\_ns std \\
\midrule
v14        & 6-fold & \bb{0.5273} & 0.570 & 0.34 \\
v14        & fold6  & 0.4824 & 0.121 & -- \\
\midrule
v15\_btc   & 6-fold & 0.5171 & 0.498 & 0.18 \\
v15\_btc   & fold6  & 0.4885 & 0.305 & -- \\
\midrule
v15\_multi & 6-fold & 0.5120 & 0.616 & \bb{0.14} \\
v15\_multi & fold6  & 0.5145 & 0.573 & -- \\
\midrule
Kronos     & fold6  & \bb{0.5345} & 0.280 & -- \\
\bottomrule
\end{tabular}
\end{table}

% ----- Table A2 -----
\begin{table}[h]\centering
\caption{Trading performance (cumulative return \%). LS = long-short, LO =
long-only with confidence gate $\theta$; ret = principal return, AER = excess
return vs BTC buy\&hold (B\&H). On fold6 the B\&H benchmark is negative
($-5.3\%$), so a positive AER means lost-less-than-holding. All models share
identical strategies and slippage.}
\label{tab:strat}
\begin{tabular}{llccccc}
\toprule
 & & LS@0.1\% & \multicolumn{2}{c}{LO@0.5\%} & Magnitude & B\&H \\
\cmidrule(lr){4-5}
Model & Period & ret & ret & AER & ret & ret \\
\midrule
v14        & 6-fold & +0.7  & +3.5 & -3.7 & +1.3  & +7.2 \\
v14        & fold6  & -13.2 & -2.7 & +2.7 & -11.0 & -5.3 \\
\midrule
v15\_btc   & 6-fold & -3.1  & \bb{+4.0} & -3.2 & -0.8  & +7.2 \\
v15\_btc   & fold6  & -7.4  & -3.2 & +2.1 & -6.7  & -5.3 \\
\midrule
v15\_multi & 6-fold & -5.7  & -2.6 & -9.8 & -4.5  & +7.2 \\
v15\_multi & fold6  & -8.5  & -8.1 & -2.7 & -8.6  & -5.3 \\
\midrule
Kronos     & fold6  & \bb{+4.7} & +0.7 & \bb{+6.1} & +4.3 & -5.3 \\
\bottomrule
\end{tabular}
\end{table}

\paragraph{Part A takeaways.}
(1) Target normalisation halves regime memorisation (Long\_ns std
$0.34\!\to\!0.18$) without multi-asset inputs.
(2) Multi-asset (ETH+XRP) does not help --- v15\_btc $>$ v15\_multi on DA and
returns.
(3) On the only clean out-of-sample window (fold6, 2025), the pretrained Kronos
baseline outperforms all our VLA variants on DA and risk-adjusted excess return.

%==========================================================================
\section*{Part B --- Macro experiments (intraday 1h, market hours, 2 folds)}
%==========================================================================

US market hours only (BTC resampled to the SPY/TLT 09:30--16:00 ET grid), 2
walk-forward folds (2024, 2025). Each row co-predicts BTC with the listed macro
ETFs; metrics are for \textbf{BTC only}. TLT/IEF/SHY are 20y / 10y / 2y US
Treasury ETFs; SPY is the S\&P500 ETF. \textbf{Caveat:} only $\sim$50 test
samples per fold $\times$ 2 folds --- low statistical power.

% ----- Table B1 -----
\begin{table}[h]\centering
\caption{Directional metrics (BTC, $k{=}6$, 2-fold avg). DA below 0.5 indicates
the intraday market-hours signal is at or below chance.}
\label{tab:macro-dir}
\begin{tabular}{lcc}
\toprule
Model & DA & Long\_ns \\
\midrule
baseline (BTC)    & 0.4646 & 0.637 \\
BTC+SPY+TLT (20y) & \bb{0.4720} & 0.486 \\
BTC+SPY+IEF (10y) & 0.3634 & 0.713 \\
BTC+SPY+SHY (2y)  & 0.4188 & 0.491 \\
\bottomrule
\end{tabular}
\end{table}

% ----- Table B2 -----
\begin{table}[h]\centering
\caption{Trading performance (BTC, cumulative \%, $k{=}6$, 2-fold avg). The B\&H
benchmark is $-5.5\%$ (BTC fell in both Dec test months), so a positive AER means
lost-less-than-holding, not a real profit --- every principal return is negative.}
\label{tab:macro-strat}
\begin{tabular}{lccccc}
\toprule
 & \multicolumn{2}{c}{LS@0.1\%} & \multicolumn{2}{c}{LS@1.0\%} & B\&H \\
\cmidrule(lr){2-3}\cmidrule(lr){4-5}
Model & ret & AER & ret & AER & ret \\
\midrule
baseline (BTC)    & -2.7 & +2.8 & -2.3 & +3.2 & -5.5 \\
BTC+SPY+TLT (20y) & -2.9 & +2.6 & -0.5 & \bb{+5.0} & -5.5 \\
BTC+SPY+IEF (10y) & -5.1 & +0.5 & -3.4 & +2.1 & -5.5 \\
BTC+SPY+SHY (2y)  & -4.7 & +0.8 & -2.4 & +3.1 & -5.5 \\
\bottomrule
\end{tabular}
\end{table}

\paragraph{Part B takeaways.}
(1) Adding bonds/S\&P does \emph{not} improve BTC prediction; baseline ($0.465$)
ties or beats the macro variants, and IEF (10y) is worst ($0.363$).
(2) All DAs are $<0.5$ and all principal returns are negative; positive AER is
only an artefact of a falling benchmark.
(3) With $\sim$100 test samples total, these results are indicative at best --- a
15m/30m re-run is needed for statistical power.

\end{document}
